% Part: proof-theory
% Chapter: sequent-calculus
% Section: interpretation-rules

\documentclass[../../../include/open-logic-section]{subfiles}

\begin{document}

\olfileid{pt}{seq}{irl}

\olsection{د قواعدو تعبير}

د سېکوېنټ تعبير راياد کړئ: سېکوېنټ $\Gamma \Sequent
\Delta$ په !!a{structure}~$\Struct{M}$ کښې رښتينے دے، که يا لږ
تر لږه يو $!A \in \Gamma$ دروغ وي، يا لږ تر لږه يو~$!A \in \Delta$
رښتينے وي۔ د~\Log{G3c} منطقي قاعدې داسې لوستلے شو چې د خپلو
اصلي !!{formula}s د رښتينولۍ شرطونه په قاعدو کښې راټولوي۔
\RightR{\lif} ته وګورئ۔ اصلي !!{formula} ئې $!A \lif !B$ دے۔
دا رښتينے دے، که يا $!A$ دروغ وي، يا~$!B$ رښتينے وي۔ نو
سېکوېنټ $\ \Sequent !A \lif !B$ رښتينے دے، که او يوازې که
$!A \Sequent !B$ رښتينے وي۔ که د سياق !!{formula}s
$\Gamma$ او $\Delta$ د دغو سېکوېنټونو کيڼو او ښيو خواوو ته
ورزيات کړو، نو کټ مټ د \RightR{\lif} قاعدې پايله او مقدمه
ترلاسه کوو۔ د~\LeftR{\lif} صورت هم دې ته ورته دے۔ $!A
\lif !B$ دروغ دے، که $!A$~رښتينے او $!B$~دروغ وي۔ نو
$!A \lif !B
\Sequent \ $ رښتينے دے، که او يوازې که دواړه $\ \Sequent !A$
او $!B \Sequent \ $ رښتيني وي۔ د \LeftR{\lif} پايله د خپلو
مقدماتو له عطفي ترکيب سره معادله ده۔ دا بڼه ډاډمنوي چې
د~\Log{G3c} قاعدې سمې دي: که ټول مقدمات رښتيني وي، پايله
رښتينې ده۔ دا قاعدې بشپړتيا هم ډاډمنوي۔

په حقيقت کښې، د هر هغه تړونکي د رښتينولۍ شرطونه، چې ارزښت
ئې يوازې د خپلو غړو د رښتينولۍ په ارزښتونو پورې تړلے وي،
په همدې ډول د سېکوېنټونو د يوې مجموعې په توګه بيانېدې شي۔
دا له دې حقيقت څخه راځي چې په کلاسيکي منطق کښې د رښتينولۍ
هره تابع په عطفي عادي بڼه د !!a{formula} په وسيله بيانېدې شي:
د اتومي او نفي شويو اتومي !!{formula}s د فصلي ترکيبونو يو
عطفي ترکيب۔ د بېلګې په توګه، $!A
\oplus !B$ هغه وخت رښتينے وګڼئ چې د $!A$ او $!B$ څخه کټ مټ
يو رښتينے وي، يعنې «انحصاري يا»۔ نو $!A \oplus !B$ رښتينے دے،
که او يوازې که $(!A \lor !B) \land
(\lnot !A \lor \lnot !B)$ رښتينے وي؛ او دروغ دے، که
$(\lnot !A \lor !B)
\land (!A \lor \lnot !B)$ رښتينے وي۔ نو د $\oplus$ دپاره د
سېکوېنټ د حساب قاعدې داسې معرفي کولے شو:
\begin{align*}
&
\Axiom$\Gamma \fCenter \Delta, !A, !B$
\Axiom$!A, !B, \Gamma \fCenter \Delta$
\RightLabel{\RightR{\oplus}}
\BinaryInf$\Gamma \fCenter \Delta, !A \oplus !B$
\DisplayProof
\\
& \Axiom$!A, \Gamma \fCenter \Delta, !B$
\Axiom$!B, \Gamma \fCenter \Delta, !A$
\RightLabel{\LeftR{\oplus}}
\BinaryInf$!A \oplus !B, \Gamma \fCenter \Delta$
\DisplayProof
\end{align*}
او دا قاعدې به هم سمې او بشپړې وي۔ پايلې ئې رښتينې دي،
که او يوازې که ټول مقدمات رښتيني وي۔
\textbf{د سرچينې سمون (OLCMP-226):} د دويمې قاعدې
پېژندنه په غلطۍ ښۍ وه؛ پايله ئې په کيڼ اړخ کښې ده،
نو کيڼه نښه سمه شوه۔ دواړه مقدمې او پايله عين دي۔

\begin{prob}
فرض کړئ $!A \downarrow !B$ رښتينے دے، که او يوازې که نۀ
$!A$ او نۀ~$!B$ رښتينے وي («نۀ دا او نۀ هغه»، يعنې NOR)۔
د دې د \Log{G3c} قاعدې به څۀ وي؟ (داسې دوه سېکوېنټونه ومومئ
چې دواړه په يو وخت رښتيني وي، که او يوازې که $!A \downarrow !B$
رښتينے وي۔ له دې څخه د \RightR{\downarrow} قاعده ترلاسه کېږي۔
بيا داسې يو سېکوېنټ ومومئ چې رښتينے وي، که او يوازې که
$!A \downarrow !B$ دروغ وي۔ له دې څخه د \LeftR{\downarrow}
قاعده ترلاسه کېږي۔)
\end{prob}

په \Log{G3c} کښې هر تړونکے او کميت‌ښودونکے کټ مټ دوه قاعدې
لري: يوه کيڼه قاعده او يوه ښۍ قاعده۔ خو په \Log{G1c} کښې
دوه \LeftR{\land} او دوه \RightR{\lor} قاعدې دي۔ لامل دا دے
چې د ګېنټسن اصلي نظام~\Log{LK}، چې \Log{G1c} د هغې په بڼه
جوړ شوے دے، د يو مهم غير کلاسيکي منطق، يعنې شهودي منطق،
دپاره بله بڼه~\Log{LJ} لرله۔ د شهودي منطق دپاره د يو سم او
بشپړ نظام د ترلاسه کولو په موخه، په~\Log{LJ} کښې د هر
سېکوېنټ ښۍ خوا په دې محدوده ده چې له يوې !!{formula} څخه
زياتې ونۀ لري۔ په دې سره د~\Log{G3c} د \RightR{\lor} قاعدې
کارول ناشوني کېږي، ځکه مقدمه ئې په ښۍ خوا کښې دواړه $!A$
او~$!B$ لري۔ نو ګېنټسن د \Log{LJ} دپاره د \RightR{\lor}
قواعدو يوه جوړه وکاروله، او هماغه ئې د~\Log{LK} دپاره هم
وکاروله۔ شهودي بڼه~\Log{G1i} د~\Log{G1c} داسې بڼه ده
چې ټولې ښۍ خواوې ئې تر يوې !!{formula} پورې محدودې دي،
خو د شرط د کيڼې قاعدې لومړۍ مقدمه هم بدله شوې ده۔
(\Log{G3c} هم يوه شهودي بڼه لري، خو د هغې ځينې قاعدې
په~\Log{G3c} کښې له اړوندو قواعدو څخه توپير لري۔ د بېلګې
په توګه، هغه هم د \RightR{\lor} قواعدو يوه جوړه کاروي۔)
\textbf{د سرچينې سپيناوی (OLCMP-228):} يوازې د هر
تالي يوه فورمول ته محدودول کافي نۀ دي۔ په شهودي
نظام کښې د شرط د کيڼې قاعدې لومړۍ مقدمه هم بدلېږي:
تالي ئې يوازې د شرط مقدمه ده، د پايلې پاتې تالي سياق
ورسره نۀ راځي۔ د اصل G1i جدول همدا بدلون لري؛ ګنې
د غيرتشې پايلې لپاره د شرط کيڼه قاعده نۀ شي کارېدلے۔

دا ليدل اسانه دي چې د~\Log{G1c} جوړښتي قاعدې سمې دي۔
د بېلګې په توګه، که $\Gamma \Sequent \Delta$ په کيڼه خوا کښې
يوه دروغ !!{formula} يا په ښۍ خوا کښې يوه رښتينې !!{formula}
ولري، نو $\Gamma \Sequent
\Delta, !A$ هم هماغه يوه لري؛ دا چې $!A$ رښتينے دے که دروغ،
څۀ اغېز نۀ لري۔ نو که د~\RightR{\Weakening} مقدمه رښتينې وي،
پايله ئې هم رښتينې ده۔ په ياد ولرئ چې دلته معکوسه خبره سمه
نۀ ده۔ ښايي پايله ځکه رښتينې وي چې $!A$ رښتينے دے، خو بيا
دا تضمين نۀ شته چې مقدمه به هم رښتينې وي۔
\textbf{د سرچينې سمون (OLCMP-227):} د سېکوېنټ د
رښتينولۍ دواړه اړخونه په اصل کښې سرچپه وو؛ دروغ
مقدمه په کيڼ او رښتينې پايله په ښي اړخ کښې وي۔ د ښي
کمزوري کولو قاعده سمه ده؛ يوازې د هغې معنايي دليل
د دې برخې له لومړني تعريف سره برابر شو۔

جوړښتي قاعدې ولې ولرو؟ انقباض (\RightR{\Contraction}،
\LeftR{\Contraction}) د بېلګې په توګه د $\ \Sequent !A
\lor \lnot !A$ د !!{prove} کولو دپاره اړين دے۔ که له
$!A \Sequent !A$ څخه پيل وکړئ، لومړے د~\RightR{\lnot} په وسيله
$\Sequent !A, \lnot !A$ ترلاسه کوئ۔ بيا د \RightR{\lor} دواړه
بڼې پرله‌پسې کارولے شئ: يو ځل د $!A \lor \lnot
!A$ د ترلاسه کولو دپاره له $\lnot !A$ څخه، او بل ځل د
$!A \lor \lnot !A$ د ترلاسه کولو دپاره له~$!A$ څخه۔ په دې توګه
$\ \Sequent !A \lor \lnot !A, !A \lor \lnot !A$ ترلاسه کېږي۔
نو په~\Log{G1c} کښې د $\Sequent !A \lor \lnot !A$ د ترلاسه
کولو دپاره بايد د يوې !!{formula} دوه څرګندېدنې په !!a{proof}
کښې «منقبض» کولے شئ:
\[
\Axiom$!A \fCenter !A$
\RightLabel{\RightR{\lnot}}
\UnaryInf$\fCenter !A, \lnot !A$
\RightLabel{\RightR{\lor}}
\UnaryInf$\fCenter !A, !A \lor \lnot !A$
\RightLabel{\RightR{\lor}}
\UnaryInf$\fCenter !A \lor \lnot !A, !A \lor \lnot !A$
\RightLabel{\RightR{\Contraction}}
\UnaryInf$\fCenter !A \lor \lnot !A$
\DisplayProof
\]
په حقيقت کښې، داسې معتبر سېکوېنټونه شته چې له کمزوري کولو
يا انقباض پرته ئې په~\Log{G1c} کښې !!{prove} کولے نۀ شئ۔
د بېلګې په توګه، په~\Log{G1c} کښې د سېکوېنټ
$!A \Sequent !B \lif !A$ !!a{proof} بايد په \RightR{\lif}
پاے ته ورسېږي (ځکه يوازې $!B \lif !A$ اصلي کېدے شي)، او
دا د مقدمې په توګه $!B, !A \Sequent !A$ غواړي۔ دا بيا د
\LeftR{\Weakening} غوښتنه کوي، څو ئې له بديهي اصل
$!A \Sequent !A$ څخه ترلاسه کړو:
\[
\Axiom$!A \fCenter !A$
\RightLabel{\LeftR{\Weakening}}
\UnaryInf$!B, !A \fCenter !A$
\RightLabel{\RightR{\lif}}
\UnaryInf$!A \fCenter !B \lif !A$
\DisplayProof
\]

په \Log{G3c} کښې د انقباض او کمزوري کولو قواعدو ته هېڅ
اړتيا نۀ شته۔ کمزوري کول په بديهي اصولو کښې له مخکې شامل
دي۔ د انقباض له کارولو څخه په زياتره صورتونو کښې داسې
ډډه کولے شو چې په دوو مقدمو لرونکو قواعدو کښې د سياق $\Gamma$ او $\Delta$
دوه نقلونه په يوه کښې يوځاے کړو۔ دواړه \Log{G1c} او \Log{G3c}
د «سياق شريکولو» داسې قاعدې لري۔ په يوه مقدمه لرونکو
قواعدو کښې انقباض يوازې د \RightR{\lor}، \LeftR{\land}، \RightR{\lexists}
او~\LeftR{\lforall} دپاره اړين دے۔ په~\Log{G3c} کښې د دې
قواعدو بڼې انقباض په بشپړه توګه بې‌اړتيا کوي۔ يو نظام
\Log{G2c} شته (\olref{tab:G2c} وګورئ) چې دوې مقدمې لرونکې
قاعدې ئې بېل بېل سياقونه لري۔ په \Log{G2c} کښې انقباض
له~\Log{G1c} څخه لا زيات مهم دے۔

\subfile{rules-G2c}

\end{document}
